\documentclass[conference]{IEEEtran}
\IEEEoverridecommandlockouts
% The preceding line is only needed to identify funding in the first footnote. If that is unneeded, please comment it out.
\usepackage{cite}
\usepackage{amsmath,amssymb,amsfonts}
\usepackage{algorithmic}
\usepackage{graphicx}
\usepackage{textcomp}
\usepackage{xcolor}
\usepackage{booktabs}   % \toprule \midrule \bottomrule
\usepackage{array}      % p{} columns behave properly
\usepackage{url}        % \url{} in the bibliography

%  Float placement 
% Without these, LaTeX queues tables up and dumps them at the end of the
% document. They let it put more floats per page and accept pages that are
% mostly float.
\setcounter{topnumber}{3}            % up to 3 floats at the top of a page
\setcounter{bottomnumber}{2}
\setcounter{totalnumber}{4}          % up to 4 floats on one page
\renewcommand{\topfraction}{0.9}     % a float may take 90% of the page top
\renewcommand{\bottomfraction}{0.8}
\renewcommand{\textfraction}{0.07}   % a page needs only 7% real text
\renewcommand{\floatpagefraction}{0.7}
\renewcommand{\dbltopfraction}{0.9}       % same, for full-width (table*) floats
\renewcommand{\dblfloatpagefraction}{0.7}
% 
\def\BibTeX{{\rm B\kern-.05em{\sc i\kern-.025em b}\kern-.08em
    T\kern-.1667em\lower.7ex\hbox{E}\kern-.125emX}}



\begin{document}

\title{The state of ML in Pokémon: the benefits of a bring predictor\\
}

\author{\IEEEauthorblockN{José Ignacio Madriz Ortiz}
\IEEEauthorblockN{Sergio Andrés Zúñiga Castillo}\\
\text{Instituto Tecnológico de Costa Rica}\\
\text{Escuela de Computación}

}


\maketitle


\begin{abstract}
\textbf{Every competitive Pok\'emon VGC match opens with Team Preview, where both
players see each other's six-Pok\'emon teams and simultaneously commit four of
their own, two of which will lead. The choice is small enough to enumerate: 90 alternatives per player, yet it constrains every turn that follows. This article reviews fourteen works on machine learning for competitive Pok\'emon
and adjacent domains and finds the decision effectively unmeasured: most agents
run in formats where it does not exist, the one benchmark that models it folds
it into the battle policy's action space without isolating its contribution,
and the only system reporting a Team Preview component publishes no result for
it. We organise the literature along three axes: the problem domain,
reinforcement-learning technique, and the available environments and data, and build a methodology from it. The bring is formalised as a two-step finite-horizon Markov Decision Process over the 90 alternatives with a sparse
terminal reward, to be trained offline on human battle logs. Direct measurement
of one Regulation~M-A log set finds 84.1\% of recorded brings complete, with
disclosure correlated to the outcome, so supervision is posed as a partial-label
problem instead of being filtered. Success is defined as clearing the 55\% win
rate at which the balance literature classifies a strategy as dominant.}

\end{abstract}

\begin{IEEEkeywords}
Pokémon VGC, Reinforcement Learning, Artificial Intelligence, Meta game, Pokémon Bring

\end{IEEEkeywords}


\section{Introduction}
In recent years, Artificial Intelligence has become an increasingly relevant subject of study, leading to many branching investigations across different areas ranging from medical assistants, evaluation tools, every-day tasks, and, of course, games. The many different genres and styles across board games, video games, and other game theory problems lead to a variety of AI solutions among these \cite{hu2024}. One of these popular video games in question is Pokémon. \textbf{Pokémon} is a series of creature-collector role-playing games centered around the \textbf{turn-based battles} that such creatures participate in. For a long time, a competitive scene has existed for this franchise, called Pokémon Video Game Championship (VGC), which has gained increasing popularity in recent years. With such a rise, and thanks to its mechanically simple turn-based combat but ever increasing meta-game, it has become a popular subject in the AI for gaming community \cite{hu2024}. \\

A Pokémon battle in the official VGC format consists of 2 teams, both with 6 Pokémon, but only 4 of which are sent into each battle.  Every Pokémon has 1 or 2 types, 4 moves, 1 ability, 1 item, a distribution of points among 6 statistics, and a nature (also called stat alignments) \cite{vgchandbook}. This alone, depending on the version of the format, leads to $\sim10^{139}$ (or $\sim10^{215}$ in bigger formats) \cite{karten2026} possible team combinations per player. That, plus the possible per-turn branching ( $\sim10^{12}$ possible outcomes),  $\sim10^{622}$ possible battle states, and a continuously evolving metagame, make Pokémon VGC a non-trivial domain worth analyzing \cite{angliss2026,karten2026} .\\ 

The present research aims towards analyzing and creating an AI algorithm capable of predicting the \textbf{bring} (the 4 Pokémon bring, including the first 2 to lead) of the enemy team and answering with one of the best possible brings. This seemingly simple process's complexity was best explained by Reis et al. as a "system governed by interdependence task model: knowing how to play is needed for good unit selection, as knowing what units are the most probable to be used affects the game itself" \cite{reis2021}; creating a more complicated Rock-Paper-Scissors situation, one bring may defeat another and so on. This sub-problem of a Pokémon battle is part of the arguably most important turn in an average match and drastically affects the outcome of a battle. Thankfully, it is explicit and countable, presenting \textbf{90 possible bring decisions per player} ($C(6,2) * C(4,2)$) \cite{angliss2026}.\\

AI applications for Pokémon VGC have progressed in 4 distinct ways: 1. Early work showed that reinforcement learning could acquire core battle concepts such as type advantage and beneficial switching without prior domain knowledge \cite{simoes2020}. 2. Self-play policy optimization later reached a Glicko-1 rating (rating of a user's performance on an online ladder) of 1677, approximately 72\% GXE (chance of winning against a random opponent from the ladder) against the best humans' 74\% to 90\% GXE\cite{grigsby2025}, at modest computational cost \cite{huang2019}. 3. More recently, offline reinforcement learning on large sequence models achieved 64–80\% GXE against human opponents, placing the agent within the top 10\%  of active players without any explicit search \cite{grigsby2025}, 4. While search-based agents combining Monte Carlo Tree Search with opponent set prediction exceeded 80\% GXE in Generation 9 OU (a format featuring over used and centralizing Pokémon) \cite{karten2026}. Given this maturity, we treat the battle policy as a reusable component rather than a contribution from this research, following \cite{reis2024} and \cite{reis2025}, which both fix an existing battle agent in order to study the surrounding decision problems. Although, the mentioned solutions are for the singles battle format, not the official VGC rules which are double battles. For the VGC rules, the best approach found was VGC-Bench as a way to measure our own model's success. \\

What has \textbf{not been solved yet} is a reliable way of determining the enemy bring and one of the best possible answers. As Table~\ref{tab:bring} shows, in most of the reviewed systems the decision does not arise at all: the environment provides two Pok\'emon per side
\cite{simoes2020}, or the format is singles, where all six Pok\'emon are
available in every battle \cite{huang2019, grigsby2025, karten2026}. The only
system that operates on the decision is VGC-Bench, which offers two
alternatives selectable by configuration: substituting a random bring, or
unifying the opening choice into the battle policy's action space as two joint
switch-in actions \cite{angliss2026}. Neither alternative predicts the
opponent's bring nor optimizes a response to it, and neither is accompanied by
a metric isolating the contribution of that choice. The single exception among the works reviewed here is Rodriguez et al. \cite{rodriguez2024}, whose decision-support tool reports a Team Preview
recommendation produced by what the authors describe as ``the comprehensive
genetic algorithm for the team preview''. That phrase appears exactly once, in
a list enumerating the output tags of the tool's analysis module. The algorithm
itself is never specified: the paper provides no representation, no fitness
function, no training procedure and no evaluation of that component, and
reports no result attributable to it. To the best of our knowledge this is the
entire published record of an attempt to optimize the bring decision
deliberately, rather than to absorb it into a battle policy. The gap is
therefore not one of interest but of measurement: the decision is acknowledged
as consequential \cite{reis2021, angliss2026}, it is made by deployed agents
\cite{angliss2026}, and it has never been isolated, optimized as an objective
in its own right, or reported against any metric.\\
This document revises the current state of the art in the domains adjacent to the problem, ways to approach it with reinforcement learning, the datasets to use and a detailed explanation on the proposed methodology for the project. 

\subsection{Research Question}
Can a reinforcement-learning policy, given both full 6-Pokémon teams at Team Preview in a VGC-adjacent fixed-team format, select a bring (4 of 6 with a designated lead pair) that achieves a win rate of at least 55\% against a population-usage baseline, over 1000 battles per team pair played in both directions with an identical fixed battle policy on both sides? 
\subsection{Hypothesis}
A reinforcement-learning policy trained on human Team Preview decisions will select brings that win at least 55\% of battles against a population-usage baseline under an identical fixed battle policy. 
\subsection{Objectives}
\begin{itemize}
\item Formalize the Team Preview decision as a Markov Decision Process over
the 90 legal brings, specifying the state space, the action space, the
transition dynamics and a justified reward function.
\item Implement a reinforcement-learning policy that, given both team sheets,
estimates the opponent's likely bring and selects a response from the 90
alternatives.
\item Define and justify an evaluation protocol that captures both the
immediate matchup advantage and the final battle outcome, and apply it to
compare the policy against a population-usage baseline at the 55\% threshold.
\end{itemize}

\subsection{Literature Gaps and Value of our Investigation}

Pok\'emon VGC, as an investigation environment, can be tackled from many
distinct angles within AI fields \cite{hu2024}. This is proven by the many
papers investigating possible implementations on the different parts of the
game, such as battling \cite{simoes2020, huang2019, grigsby2025, karten2026} or
team-building \cite{reis2024}. To effectively create developments in these
areas, the environments have to be reduced to a manageable search space, which
leads to other aspects of the game being left out by researchers: simplified
simulators isolate a single learning objective \cite{simoes2020}, team building
is declared out of scope \cite{angliss2026}, and the battle policy is frozen so
that the surrounding decisions can be studied \cite{reis2024, reis2025}.

Even if team-building and battling are arguably the two most important pillars
of VGC, there is one middle step that can have a significant impact on the
outcome of a specific match: analyzing your enemy's team and deciding your
bring. Reis et al. place precisely this step between team building and the
battle itself, as a dedicated Selection Phase \cite{reis2021}. This process as a
whole is often referred to as \textbf{Team Preview}, and occurs just before a
match begins \cite{vgchandbook}. A bad guess on Team Preview can bring a whole
competitive strategy to the ground, undermining the value of battling and
team-building; Angliss et al. observe that this single decision ``can
drastically alter the dynamics of the battle'' \cite{angliss2026}.

Despite this, the step is routinely set aside. Many battling agents operate in
formats where it does not exist at all \cite{simoes2020, huang2019,
grigsby2025}, and the one benchmark that does model it offers a configuration
that replaces the bring with a random selection in order to encourage
exploration during training \cite{angliss2026}. Whether battle skill can
compensate for a poor bring has never been measured. Nevertheless, in the
process of creating a complete AI environment for the scene, it is important to
have mastered every aspect of it, particularly given that selection and
battling are mutually dependent \cite{reis2021}.

The value of our investigation is inherent from this: the crucial opening
decision, informally known among players as ``Turn 0'', has not been isolated,
optimized as an objective in its own right, or reported against any metric. Our
aim is to create an algorithm capable of analyzing the meta-game and both
teams' strengths and weaknesses, and coming up with a suitable answer based on
that reasoning. If the investigation reaches a desirable outcome, it could also
contribute to more robust battling or team-building algorithms; Angliss et al.
note that a strong VGC agent could supply a reward signal for searching the team
configuration space \cite{angliss2026}.

\emph{Quick note on scope:} The following fall explicitly outside the scope of this work
\begin{itemize}
\item {Team building.} Constructing six Pok\'emon from a space of
approximately $10^{139}$ configurations \cite{angliss2026} is left open, as it
is in \cite{angliss2026} and \cite{reis2021}. Teams are treated as given.
\item{Training a battle policy.} Existing agents are reused
\cite{grigsby2025, angliss2026, karten2026}, following the same design choice
made in \cite{reis2024} and \cite{reis2025}.
\item{Metagame and rules balance} \cite{reis2021, reis2024, reis2025, hu2024}.
\item{Large language model approaches}, which \cite{karten2026} and
\cite{grigsby2025} report as uncompetitive over long horizons in this domain.
\item{Superhuman play.} Following \cite{li2025} and \cite{rodriguez2024},
the target is decision support for a human player.
\end{itemize}
\section{Related works}

The reviewed literature is organized along three axes: the problem domain,
covering competitive Pok\'emon and specifically the selection of a bring at
Team Preview; the techniques associated with a reinforcement-learning approach;
and the datasets and simulation environments available for this domain.
Following the survey methodology described by Hu et al. \cite{hu2024}, works
were collected from Google Scholar and the various tools provided by Instituto Tecnológico de Costa Rica's digital library using the search terms such as "Pokémon VGC", "Reinforcement Learning in video games", "meta-game analysis" or "Pokémon AI battles", additionally, the AI tool Consensus helped in finding certain related papers under different keywords. Research was restricted to works published between 2019 and 2026, with the exception of one paper \cite{galway2008}, and complemented by backward citation chasing from the VGC AI Competition literature and by direct consultation with a competitive VGC practitioner active in AI research. Works were included when they addressed Pok\'emon as an application domain, a reinforcement-learning or evolutionary technique applicable to sequential decision-making under partial information, or a simulation environment or dataset for competitive Pok\'emon;
meanwhile others were excluded when they provided no relevant new information and had older publication dates. Table~\ref{tab:related} summarizes the resulting
set along the dimensions required for comparison, and Table~\ref{tab:bring}
details how each reviewed system treats the Team Preview decision in
particular. The three axis of this investigation are developed in the upcoming subsections.

% Include: Reference, year, technique, enviorments, main metric, reported result, identified limitation
To better understand the scope and relevance of each of the consulted papers, Table-\ref{tab:related} compares them in categories based on what they did, where they did it, how they measured the results, what the results were and limitations regarding further investigation.
\begin{table*}[!t]
\centering
\caption{Comparison of related work on machine learning for competitive Pok\'emon and adjacent domains.}
\label{tab:related}
\scriptsize
\setlength{\tabcolsep}{3pt}
\renewcommand{\arraystretch}{1.15}
\begin{tabular}{@{}p{1.4cm}p{0.5cm}p{3.0cm}p{2.8cm}p{2.0cm}p{3.5cm}p{3.3cm}@{}}
\toprule
\textbf{Ref.} & \textbf{Year} & \textbf{Technique} & \textbf{Dataset / environment} & \textbf{Main metric} & \textbf{Reported result} & \textbf{Identified limitation} \\
\midrule

Galway \emph{et al.} \cite{galway2008} & 2008 &
Survey of neural networks, evolutionary computation and RL for game-agent control &
Academic digital-game research on games before 2008 &
Survey, no metric &
Neural network and evolutionary hybridization predominates. Stochastic games give noisy fitness. Pre-initialized populations yield more robust controllers &
17 years old, predates deep RL, no Pok\'emon mention \\
\addlinespace

Huang and Lee \cite{huang2019} & 2019 &
Proximal Policy Optimization with Generalized Advantage Estimation, self-play; actor and critic with 1.33M parameters &
Pok\'emon Showdown \texttt{gen7randombattle}; 3-team fixed metagame &
Win/loss vs.\ agents; Glicko-1 on the public ladder &
612--388 vs.\ a tree-search agent; Glicko-1 1677 ($\approx$72\% GXE); fixed-team specialist wins all 9 matchups (0.555--0.99) but falls to 77/500 in the general format &
Singles, Gen 7; only 3 hand-picked teams; no Team Preview; shaping constants not ablated \\
\addlinespace

Sim\~oes \emph{et al.} \cite{simoes2020} & 2020 &
Deep competitive RL (WPL$_\theta$, GIGA$_\theta$), self-play over 2M episodes; $R = R_d + R_f - R_t$ &
Pok\'emon Battle Environment (PBE), Gym-compatible; 1 active + 1 benched Pokémon &
Win rate in a fixed disadvantageous scenario &
GIGA$_\theta$ 100\% win rate; converges faster than WPL$_\theta$; with all 18 types WPL$_\theta$ fails while GIGA$_\theta$ matches the expert baseline &
No weather, status, field effects or items; 2 Pok\'emon per side; one hand-built scenario; no bring selection \\
\addlinespace

Reis \emph{et al.} \cite{reis2021} & 2021 &
Competition framework (VGC AI Framework; Design Constraint Language) &
VGC AI engine (151 templates, $\geq$165 moves) &
No experiment mentioned &
Defines three tracks and a four-module pipeline including a Selection Phase and a Team Predictor &
No baseline published for any track; simplified non-Showdown engine \\
\addlinespace

Hu \emph{et al.} \cite{hu2024} & 2024 &
Survey, Systematic review and taxonomy (2016 - 2024) &
40 games, 20 platforms, 26 design environments &
Survey, no metric &
Pok\'emon VGC listed as an AI testbed and as one of three platforms for game tuning and balance &
One paragraph on Pok\'emon, no metrics; predates Metamon, VGC-Bench and Pok\'eAgent \\
\addlinespace

Li and Wellman \cite{liwellman2024} & 2024 &
Meta-game evaluation via bootstrapped empirical games; Gumbel Information-Set Monte-Carlo Tree Search &
Deal-or-No-Deal negotiation; 17 MARL algorithms $\times$ 10 seeds &
NE-Regret, uniform score, NE-NBS, best-response graphs & The best method combines search during both training and evaluation. Agents trained by self-play show signs of having over-fitted to themselves &
Small negotiation domain, not Pok\'emon; requires a full cross-play payoff matrix; no evolutionary method evaluated \\
\addlinespace

Reis \emph{et al.} \cite{reis2024} & 2024 &
Two agents compete: one builds teams to maximise win rate, the other retunes Pok\'emon attributes to restore variety. Builders use game theory and evolutionary search &
VGC AI Framework; 3-Pok\'emon teams; 1M simulated matches; fixed Minimax battle policy &
Win-prediction accuracy; agent ranking &
The win predictor reaches 85.4\% accuracy on unseen team pairs. The strongest builder, Meta Counter, beats every other builder and wins the simulated championship &
Synthetic simplified environment, no human data; team building rather than the bring decision; results vary widely between runs \\
\bottomrule
\end{tabular}
\end{table*}

\begin{table*}[!t]
\centering
% Shares the number of the table above: \addtocounter cancels the step
% that \caption performs, so this prints as "TABLE I (continued)" and the
% next float in the document is still TABLE II. No \label here on purpose.
\addtocounter{table}{-1}
\caption{Comparison of related work on machine learning for competitive Pok\'emon and adjacent domains (continued).}
\scriptsize
\setlength{\tabcolsep}{3pt}
\renewcommand{\arraystretch}{1.15}
\begin{tabular}{@{}p{1.4cm}p{0.5cm}p{3.0cm}p{2.8cm}p{2.0cm}p{3.5cm}p{3.3cm}@{}}
\toprule
\textbf{Ref.} & \textbf{Year} & \textbf{Technique} & \textbf{Dataset / environment} & \textbf{Main metric} & \textbf{Reported result} & \textbf{Identified limitation} \\
\midrule

Rodriguez \emph{et al.} \cite{rodriguez2024} & 2024 &
Compares two ways of training the same agent: Double Deep Q-Learning against an evolutionary algorithm that evolves the network itself (NEAT). Also creates a recommend-move tool to a player &
Pok\'emon Showdown, VGC 2022 doubles; 8 teams from the 2022 World Championship &
Accumulated reward; wins and losses per opponent team &
After 5 days NEAT scores slightly higher than DDQN. DDQN wins 53 of 80 battles; NEAT wins 204 of 480 &
No baselines, seeds or confidence intervals; reward never formally specified; comparison acknowledged as unfair; only 10 battles per team \\
\addlinespace

Grigsby \emph{et al.} \cite{grigsby2025} & 2025 &
Large Transformer models (up to 200M parameters) trained human battles, improved with self-play data. Uses memory of the current battle, no simulating future turns &
Showdown Gen 1--4 OU; 475k reconstructed battles, 950k trajectories, 38M steps &
GXE (chance of beating a random ladder opponent); Glicko-1 &
64--80\% GXE against human players, placing it in the top 10\% of the ladder; peaked at rank \#31 globally in Gen 1 OU &
Singles only, no doubles or VGC; no Team Preview by design; no team building or bring decision \\
\addlinespace

Reis \emph{et al.} \cite{reis2025} & 2025 &
A genetic algorithm tunes the game's rule parameters. Each candidate set of rules is scored by having a fixed tree-search agent play battles under them &
VGC AI Competition 2.0; 3-Pok\'emon teams; generated team pairs &
How close move usage gets to a target rate &
Converges in 7--25 generations. The tuned rules score 0.88--0.98 against 0.72--0.94 for the official human-set parameters &
Optimises rules, not selection; win rate left out of the fitness, producing extreme values (authors' caveat) \\
\addlinespace

Li \emph{et al.} \cite{li2025} & 2025 &
Review of multi-agent RL in video games; proposes five dimensions for rating how hard a game is to learn &
Sports, FPS, RTS and MOBA games (StarCraft~II, Dota~2, GRF) &
Survey, no metric &
Rates difficulty by observability, state space, action space, reward sparsity and number of agents. Documents the gap between academic RL and shipped game AI &
No Pok\'emon and no turn-based simultaneous-move games; nothing on pre-game selection or drafting \\
\addlinespace

Oblander \cite{oblander2025} & 2025 &
Neural network that maps each 6-Pok\'emon team to a vector, then predicts the winner from the two vectors. Corrects for player skill &
Showdown ranked VGC; 11\,738\,418 matches, 400\,009 players, 702 Pok\'emon &
Prediction accuracy; AUC &
66.43\% accuracy predicting the winner from both teams, against 65.54\% using player skill alone. Players base 59.3\% of their team choice on what is popular &
Models the 6-Pok\'emon team, never the 4 brought; unordered sets \\
\addlinespace

Angliss \emph{et al.} \cite{angliss2026} & 2026 &
Agents trained by copying human play, then improved by playing against themselves and against a growing pool of opponents &
VGC-Bench: Showdown VGC doubles, Open Team Sheets, PettingZoo; 700k+ human battle logs &
Win rate against each other; exploitability; performance on unseen teams &
An agent trained on one team beats agents trained on many (about 70\% win rate) but generalises worse. A purpose-built counter-strategy beats any of them almost always. Beat an intermediate human 5/5 and an advanced human 2/5 &
Team building out of scope; near-total exploitability; Open Team Sheets imply near-full information; only 5 games per human evaluation \\
\addlinespace

Karten \emph{et al.} \cite{karten2026} & 2026 &
Benchmark comparing heuristic agents, RL agents, language-model agents and tree search on a shared server &
Pok\'eAgent Challenge (Showdown Gen 1 and Gen 9 OU); 4M human and 18M self-play trajectories; 200k+ teams &
GXE; Glicko-1; a full-history rating that is less noisy for small agent pools &
The tree-search agent Foul Play exceeds 80\% GXE in Gen 9 OU and won its track 50--14. RL agents beat language-model agents &
Singles-focused; no standalone bring or lead metric; ratings noisy without large samples; metagame drift \\
\bottomrule
\end{tabular}
\end{table*}

\subsection{Problem and domain}
The A axis refers to the scope of the problem and the domains it is part of. How has the problem of selecting a suitable bring to a Pokémon VGC match been handled by similar research, which domains of research are relevant for this investigation and what has been solved or what has not been investigated as in depth.  
\subsubsection{How has this problem been handled previously?}

The bring decision is a problem often overlooked, avoided, or not dealt with as a problem due to format differences or design decisions. The Table-\ref{tab:bring} shows how the previous AI models have had to deal with, if at all, with this particular problem.

%Handling of bring decision on the papers
\begin{table*}[!t]
\centering
\caption{Treatment of the Team Preview (bring) decision across the reviewed AI models or experiments.}
\label{tab:bring}
\scriptsize
\setlength{\tabcolsep}{3pt}
\renewcommand{\arraystretch}{1.15}
\begin{tabular}{@{}p{3.3cm}p{1.6cm}p{3.4cm}p{8.2cm}@{}}
\toprule
\textbf{System} & \textbf{Ref.} & \textbf{Format} & \textbf{Treatment of the bring decision} \\
\midrule
GIGA$_\theta$, WPL$_\theta$ & \cite{simoes2020} & PBE (1 active + 1 benched) & No decision taken: Only two Pokémon per team, set leads.\\
\addlinespace
RL-rb, RL-meta & \cite{huang2019} & Showdown Gen 7  random singles; 3 fixed teams& No decision taken: in singles all six Pok\'emon are usable in every battle. Team Preview is never mentioned and lead is forced random.\\
\addlinespace
Metamon (IL, RL, SynRL) & \cite{grigsby2025} & Showdown Gen 1--4 OU singles & Absent by design: these generations predate Team Preview, and team order is fixed when the team is built. \\
\addlinespace
Foul Play & \cite{karten2026} & Showdown Gen 9 OU singles & No four-of-six choice, since singles brings all six. Performs opponent \emph{set} prediction (hidden moves and items), which is opponent modelling rather than bring prediction. \\
\addlinespace
PA-Agent & \cite{karten2026} & Showdown Gen 1 OU singles & No bring decision necessary. Picks a team from a candidate pool by tournament selection against 50+ line-ups, retaining those above a 60\% win rate.\\
\addlinespace
VGC-Bench agents (BC, BCSP, BCFP, BCDO, PPO) & \cite{angliss2026} & Showdown VGC doubles, Open Team Sheets & \textbf{Considered, but made part of the battle policy}: Team Preview is unified into the action space as two joint switch-in actions, without a distinct representation and without any reported metric isolating its contribution. A configuration toggle replaces the bring with a random selection.\\
\addlinespace
State-space analysis only & \cite{karten2026} & Gen 9 VGC & The Battling Track is singles, and VGC appears only in a complexity table.\\
\addlinespace
DDQN, NEAT, and an unnamed genetic algorithm & \cite{rodriguez2024} & Showdown VGC 2022 doubles & A genetic algorithm is reported to produce a Team Preview recommendation, but no further explanation about that specific module is given.\\
\addlinespace
Team builders (Individual Pok\'emon Counter, Minimax Builder, Meta Counter) & \cite{reis2024} & VGC AI Framework, three-Pok\'emon teams from a roster & No bring: these agents build teams from a roster. The Selection Phase is not implemented, and battling uses a fixed Minimax policy rather than a learned one.\\
\addlinespace
Tree-search agent with a PyGAD balance agent & \cite{reis2025} & VGC AI Competition 2.0 & No bring: the genetic algorithm optimises rule parameters rather than selection. \\
\addlinespace
Selection Phase (specification only) & \cite{reis2021} & VGC AI Framework & Defined as a task (predict the opponent's units, then choose the sub-team that answers them) but implements no agent and no baseline.\\
\addlinespace
Bilinear matchup model & \cite{oblander2025} & Showdown ranked VGC & Not an agent rather, it models the choice of the six Pok\'emon team as an unordered set, never the four brought.\\
\bottomrule
\end{tabular}
\end{table*}

Summarizing \ref{tab:bring}, the three relevant mentions on a Team Preview algorithm are \textbf{1. a genetic algorithm} producing a Team Preview recommendation appears in the tool of Rodriguez et al. \cite{rodriguez2024}, but only as an output of their analysis module; the component is not specified or evaluated anywhere in the paper, and no result is reported for it.\textbf{ 2.} Reis et al. \cite{reis2021} formalize this decision inside a \textbf{dedicated Selection Phase}, consisting of a Team Predictor that estimates the opponent's units and a Selection Policy that chooses a sub-team in response, and characterize it as a one-shot process. The framework provides the architecture but no implementation of either component: the paper reports no experiments, and the authors note that baseline solutions remain to be developed. \textbf{3.} The only one that reports how they actually solve or take on the problem is Angliss et al. \cite{angliss2026}, who simply treat it as \textbf{two joint switch-ins}, a smart way to re-utilize existing logic, but does not confront the problem nor its ramifications directly. 


\subsubsection{Related domains}
The Pokémon VGC bring problem studied here is a version of a pattern common to
competitive multiplayer games. Reis et al. observe that games with distinct
mechanics share a phase of pre-game preparation: assembling a deck
in card games, selecting a fighter in fighting games, or choosing
units in real-time strategy games, \cite{reis2021} where players may have partial observability on the enemy's options. Hero drafting in MOBAs is another increasingly common version of a similar case, addressed in a limited form by OpenAI Five
\cite{angliss2026}. In each case a player commits to an initial configuration before
play begins, under uncertainty about the opponent's choice, and that commitment
will affect every subsequent decision. This counter-selection aspect has been
formalized directly: Crane et al. define Team Counter-Selection Games, in which
candidate selection is shaped by a rotating metagame where each new selection
creates a counter to the previous one, and apply epistemic reasoning to
anticipate opposing line-ups and identify the strongest counter candidates,
evaluated on online Pok\'emon Go leagues \cite{crane2021}. This Rock-Paper-Scissors style of problems is part of the intersection between \textbf{metagame analysis}, which studies how strategy selection evolves in a population of players \cite{reis2024}, and
empirical \textbf{game-theoretic analysis}, which corresponds to the evaluation of strategies
whose value depends on the opponents they are measured against
\cite{liwellman2024}. Pok\'emon VGC is a demanding representative of this
family: the selection is combinatorial in a space of 90 possible distinct choices \cite{angliss2026}, the opponent's choice is simultaneous
and hidden, and the environment in which a population of strategies develops shifts over time \cite{hu2024, karten2026}.
\subsubsection{What has been solved and what is an open question?}
The two main areas on Pokémon VGC  are battling and team-building, both of which have been explored and solved at different levels of complexity and in distinct formats \cite{simoes2020,huang2019,grigsby2025,karten2026}. It is our exact problem which is still an open question: what to bring, given two teams. This has been addressed by various researchers, but it is not covered in depth, \cite{rodriguez2024, reis2021} or treated as a side effect of learning to battle \cite{angliss2026} and never tested and reported in an isolated environment by itself.  

\subsection{Reinforcement-learning approach techniques}
\label{sec:rltechniques}

The bring decision is a sequential choice made under uncertainty about a
simultaneous opponent action, with an outcome observable only after the fact.
This section reviews the families of methods that have been applied to
decisions of this shape in competitive Pok\'emon and adjacent domains, the
assumptions each one carries, and the conditions under which those assumptions
hold. Since no formal or complete research has been done on this particular aspect of the game, the purpose here is to establish, given the existing literature, why a reinforcement-learning formulation is the appropriate basis
for the problem posed, and which of its variants the available data and
computational budget permit. Not all of these were considered for our experiment, but given the survey-like nature of this work, they're all worth including.

\subsubsection{State of methods and tools: strengths, weaknesses, and assumptions}
\paragraph{Online reinforcement learning through self-play}
The earliest results in this domain come from agents that learn by playing
against themselves. Sim\~oes et al. train the competitive algorithms
WPL$_\theta$ and GIGA$_\theta$ over two million self-play episodes and show
that type advantage and the value of switching are acquired without any prior
domain knowledge \cite{simoes2020}; Huang and Lee reach a comparable result with Proximal Policy Optimization, at a reported cost of
approximately \$91 of computation \cite{huang2019}, and Rodriguez et
al. apply Double Deep Q-Learning to the VGC doubles format
\cite{rodriguez2024}. The strength of this family is that it requires no
labelled data, only an environment to interact with. Its assumption is that
such interaction is relatively cheap: all three works generate millions of episodes, but
Sim\~oes et al. report that WPL$_\theta$ fails to learn adequate policies once
the full eighteen-type environment is used, indicating sensitivity to the
scale of the interaction budget relative to the complexity of the state space
\cite{simoes2020}. A second limitation is specialization: Huang and Lee's agent,
retrained on a fixed set of three teams, wins every matchup in that setting but
falls to 77 wins in 500 games when returned to the general format
\cite{huang2019}, a pattern Angliss et al. reproduce at larger scale by finding
that agents trained on a single team configuration generalize worst across
unseen ones \cite{angliss2026}. Given the budgetary constraints of this work (\$0.00), and the limitations revealed by previous research, an additional angle other than restricted self-play must be implemented.

\paragraph{Offline reinforcement learning and imitation}
When there exists a registry of recorded play, an agent can be trained from logged
trajectories rather than from live interaction. Grigsby et al. reconstruct
475,000 human battles into trajectories and train sequence
models on them, reaching 64-80\% GXE against human opponents without any
search at decision time \cite{grigsby2025}. Angliss et al. report that across their empirical training,
the variants initialized by behaviour cloning on human play dominate those
trained from scratch \cite{angliss2026}. The assumption here 
is the availability of outcome-labelled data, which is precisely the condition
Angliss et al. establish for the VGC doubles format through their release of
over 700,000 open team sheet battle logs \cite{angliss2026} (which become the lifeblood of our fitness function). Thanks to that, the interaction cost that constrains the previous family
disappears, which is why this is the technique that most closely describes our experiment.

\emph{One issue} with the logs, it records the Team Preview information and the battle's turns. In VGC it is possible for a winner to not reveal all 4 brought members or for a loser to forfeit before all 4 of theirs are knocked out. This means not all logs carry complete 4-Pokémon brings, which is addressed in the methodology.

\paragraph{Evolutionary computation and neuroevolution}
Evolutionary methods have been applied to this class of problem both as
optimisers of network parameters and as search procedures over discrete
configurations. Galway et al. report that the hybridisation of neural networks
with evolutionary methods predominates in academic game research
\cite{galway2008}, and Rodriguez et al. provide a direct comparison within
Pok\'emon VGC, evaluating NEAT against DDQN on the same task
and finding the two comparable in performance \cite{rodriguez2024}. Applied to
search rather than to training, Reis et al. combine evolutionary search with
linear programming to construct counter-teams, scoring candidates with a
neural network trained to predict match outcomes \cite{reis2024}, and later
employ a genetic algorithm over rule parameters that converges within 7 to 25
generations \cite{reis2025}. Two constraints emerge from this literature and
inform the design proposed here. Stochastic
environments produce noisy fitness evaluations, which require a large number of
games per candidate before an estimate becomes reliable \cite{galway2008},
the condition that motivates the learned evaluator of
Section~\ref{sec:evaluator}. Additionally, excluding win
rate from their fitness function, which we will do, can drive the optimizer towards degenerate
parameter values \cite{reis2025}. These search families are more so components that combine with reinforcement learning and the design proposed
here treats them as a baseline
rather than as an eligible method.

\paragraph{Search-based methods}
The strongest battling agents currently rely on explicit search. The Monte
Carlo Tree Search engine in \cite{karten2026} exceeds 80\% GXE in
Generation 9 OU, and Reis et al. adopt a
Minimax agent as a fixed battle policy precisely because it can
operate under modified rules without retraining \cite{reis2024}. The biggest issue here is, once again, computational cost and time: the expense lies not in simulating a battle once but in evaluating ninety candidate decisions repeatedly, which is the cost a learned evaluator is
introduced to avoid. However, since our objective here is to clear a win rate percentage, a search-based policy will be useful for the significantly reduced number of simulations we will need to run.

\paragraph{Language-model agents}
Approaches built on LLMs are excluded based on evidence in the
literature rather than on principle: specialist agents outperform language-model agents over
the long horizons characteristic of VGC \cite{karten2026} \cite{grigsby2025}. Additionally, our expert consultation dissuaded us from this since the beginning.


This literature has made clear that the choice of algorithm cannot
be separated from the protocol used to assess it. Performance in multi-agent settings is meaningful only as a distribution
over random seeds and opponents, and that agents trained by self-play leave
detectable signs of having over-fitted to themselves \cite{liwellman2024}.
Angliss et al. measure the corresponding weakness directly in this domain,
finding agents to be close to fully exploitable by a purpose-built counter
strategy \cite{angliss2026}. Li et al. add a further consideration, arguing that
human-like and useful behaviour is frequently a more appropriate target for game
agents than superhuman performance \cite{li2025}, which aligns with the
decision-support framing adopted by Rodriguez et al. \cite{rodriguez2024} and
with the objectives stated in this work.


Four properties of the bring decision follow from the preceding review and
together determine the formulation developed in Section~\ref{sec:mdp}. The
decision is sequential and made under uncertainty about an opponent's
simultaneous choice, which is the structure a Markov Decision Process is
intended to describe. Its action space is small and closed at ninety
alternatives \cite{angliss2026}, so the scale that constrains self-play
approaches does not arise. In summary, outcome-labelled data exists for the format
\cite{angliss2026}, so the policy can be trained offline in the manner of
\cite{grigsby2025} without the interaction cost that the online family
requires. And the quantity of interest is a terminal win or loss, which is
already the reward signal a reinforcement-learning formulation optimises,
removing the need for the shaped objectives that \cite{simoes2020} and
\cite{grigsby2025} introduce to address credit assignment across long battles.
The evolutionary methods reviewed above remain relevant as a search procedure
over the same space, and are retained in that role.
\subsection{Dataset and environments}
% Must conclude an **objective number** with best pub lished result. Reasonable range to expect. Declare which result we wish to reach. Point out protocol differences and other metrics that might result in an unfair comparison. In its defect, look for possible comparisons from similar works
\label{sec:axisc}
The third axis concerns the environment in which the decision is actually made
and the record of play that environment produces. Two questions have to be
settled before any result of ours can be interpreted: whether a simulator
exists that poses the decision under study at all, and what the literature has
already achieved on it. The first has a narrow answer, since only one
environment qualifies. The second has an inconvenient one: nothing published
measures the quantity measured here, which changes what a literature baseline
can be for this work and is treated explicitly below.

\subsubsection{The environment and its published results}
Most of the reviewed work does not run in an environment where a bring exists.
Sim\~oes et al. use a battle environment with one active and one benched
Pok\'emon \cite{simoes2020}; Huang and Lee use Generation~7 random singles
\cite{huang2019}; Grigsby et al. use Generations~1--4 OU singles and exclude
Team Preview by design \cite{grigsby2025}; Karten et al. benchmark Generation~1
and Generation~9 OU, also singles \cite{karten2026}. The VGC AI Framework of
Reis et al. does contain a Selection Phase, but it runs on a simplified
non-Showdown engine with three-Pok\'emon teams \cite{reis2021, reis2024,
reis2025}, so the selection it models is not the four-of-six decision the
official format defines \cite{vgchandbook}.

Two works run VGC doubles on Pok\'emon Showdown itself: Rodriguez et al., over
eight teams from the 2022 World Championship \cite{rodriguez2024}, and Angliss
et al. \cite{angliss2026}. Of the two, only VGC-Bench supplies the four
properties this work requires at once: doubles under current regulations, an
Open Team Sheet setting in which both configurations are observable at Team
Preview and the state of (\ref{eq:posg}) is therefore well defined, a
representation in which the Team Preview decision is an action rather than an
assumption, and a released corpus of human play. No other environment in the
reviewed literature offers that combination, and the choice of VGC-Bench
follows from its absence elsewhere rather than from preference.

The results published on it, with the protocol that produced them, are as
follows. Under cross-play at 1000 episodes per matchup, five training seeds and
a reported error of approximately $\pm 0.03$, a behaviour-cloned agent
fine-tuned by self-play on a single team defeats the agents trained on 4, 16
and 64 teams at win rates of 0.699, 0.740 and 0.698. Evaluated instead on 72
teams held out of training, the ordering inverts: the single-team agent falls
to 0.405, 0.375 and 0.331, while the 64-team agent reaches 0.669, 0.578 and
0.564. Nearly every agent is measured as approximately fully exploitable by a
purpose-trained best response. In human testing the strongest agent won 5 of 5
against an intermediate player, 2 of 5 against an advanced player, and at least
one game against a multi-time World Championship competitor and regional
champion \cite{angliss2026}.

Two things follow. The evaluation protocol adopted in
Section~\ref{sec:protocol}, (1000 battles per matchup, sides exchanged,
results as distributions over seeds, team pairs held out) is taken from this
work rather than invented, because it is the only protocol with published
numbers on the same simulator. And the spread of those numbers, roughly 0.33 to
0.74, is the range within which differences in this environment have so far
been found meaningful. It is the empirical context for the target set out in
Section~\ref{sec:baseline}.

\subsubsection{The dataset}
The corpus released with VGC-Bench contains over 700\,000 Open Team Sheet
battle logs collected from Pok\'emon Showdown \cite{angliss2026}. This work
uses one file of it, \texttt{logs\_gen9championsvgc2026regma.json}, holding the
6\,905 battles, and therefore 13\,810 player-brings, played under the
Regulation~M-A ruleset of the Generation~9 VGC~2026 format.

Restricting to a single regulation is deliberate. A regulation is a legal
Pok\'emon pool, and therefore a metagame: usage distributions and counter
relationships are properties of one pool and do not carry across pools. Mixing
regulations would make the usage-statistics baseline of
Table~\ref{tab:baselines} an average over populations that never played each
other, and would corrupt the temporal partition of
Section~\ref{sec:protocol}, whose purpose is to separate early from late play
within one metagame rather than to separate one format from another.

Three properties of the registry constrain what may be claimed from it, and each
is a consequence of how it was collected. The first is that only Open Team
Sheet battles are retained, because only there can states and actions be
recovered by replaying the logged transitions \cite{angliss2026}. That is what
makes the data usable, but it also means the recorded population is the one
that plays under tournament-style disclosure rather than the ladder at large,
and the authors themselves list the near-full information this implies as a
limitation of their setting. The second is a rating floor: the collector reads
only logs of players above a rating threshold, so the sample is skewed toward
stronger play. For this work that skew runs in a convenient direction, since
the object of study is a competent bring rather than an average one, but it
means any statement made here about how players bring is a statement about
strong players specifically. The third is that the bring itself is not always
fully recorded, and is missing in a way that correlates with the outcome; that
property is measured and handled in Section~\ref{sec:observability}.

Size is a genuine constraint rather than a formality. Spread across a 90-way
decision, 13\,810 player-brings would average roughly 153 observations per
alternative if the decision were uniformly distributed, and it is not: usage
concentrates heavily, so the tail of rare brings is supported by very few
examples. This is a further reason the partial-label objective of
(\ref{eq:partial}) retains every record instead of discarding ambiguous ones.
At this scale the 15.9\% that would be discarded is not affordable, quite apart
from the bias that discarding them would introduce.

\subsubsection{The absence of a comparable published result}
Rodriguez et al. is the only reviewed work whose system contains a Team Preview
component at all: their analysis module emits a recommendation produced by a
genetic algorithm, printed in the tool's report beside the per-turn move
recommendation \cite{rodriguez2024}. The evaluation in that paper, however,
measures the battle agents and only the battle agents. Its results are position
tables for DDQN and NEAT across eight opponent teams at ten
interactions each, and its stated conclusion is a claim about battling. No accuracy, win
rate or degradation is reported for the team-preview component, and it is not
specified in enough detail to implement. The nearest thing in the literature
to the problem posed here therefore supplies no number.

What the rest of the field supplies are measurements of adjacent quantities.
Oblander predicts the winner from two six-Pok\'emon teams at 66.43\% accuracy
against a 65.54\% skill-only floor \cite{oblander2025}; Reis et al. predict the
winner of a simulated match at 85.4\% on unseen team pairs \cite{reis2024};
Grigsby et al. reach 64--80\% GXE as a battle policy \cite{grigsby2025}. None
of these is the bring. The first two condition on the full six and never on the
four brought, and the third is a battle policy in a format without a Team
Preview at all.

This is precisely why the review is met by
examining analogous datasets. Table~\ref{tab:baselinefigures} sets out those reference points
together with the specific difference that disqualifies each one. Our position
is that none of them constitutes a fair comparison and that reporting a result
against any of them would mislead rather than inform: the quantity, the data
and the protocol all differ. They are retained as an indication of the
magnitude of effect that is attainable in this domain, and for nothing further.

\subsubsection{The target adopted in its place}
\label{sec:baseline}
Since no published result exists for the quantity measured here, the objective
number of the research question cannot be inherited from a prior result. It is
taken instead from the threshold at which this literature declares a strategy
dominant.

Reis et al. state that threshold explicitly. Among the targets a balance agent
must satisfy, the Design Constraint Language of the VGC AI Competition sets
\texttt{MAX\_WIN\_RATE = .55}: a roster is admissible when the maximum win rate
attainable by any Pok\'emon or team is at most 55\%, the stated purpose being
to ensure there is ``no single dominant strategy'' \cite{reis2021}. The
complementary half of the definition appears in the later work, where the
fitness of a balanced metagame drives the average win-rate array toward
$w = 0.5, \ldots, 0.5$ \cite{reis2024}. Together the two fix a scale: 50\% is
what balance looks like, and 55\% is where this literature stops calling a
strategy balanced and starts calling it dominant.

This work adopts the same number with its role reversed. What a balance agent
treats as a ceiling to be pushed down, a selection policy treats as a floor to
be cleared. A bring policy that sustains more than 55\% against the
population-usage baseline (with the teams themselves fixed and exchanged
between sides, and every battle played out by a shared fixed policy) has by
this literature's own criterion produced an advantage that a metagame would
classify as dominant. That the advantage would arise from selection rather than
from roster composition is exactly the claim under test.


\begin{table}[!t]
\centering
\caption{Reference points for the literature baseline.}
\label{tab:baselinefigures}
\footnotesize
\setlength{\tabcolsep}{4pt}
\renewcommand{\arraystretch}{1.2}
\begin{tabular}{@{}p{2.5cm}p{1.5cm}p{3.4cm}@{}}
\toprule
\textbf{Reference point} & \textbf{Value} & \textbf{Why it is not a comparison} \\
\midrule
Team Preview component \cite{rodriguez2024} & none reported & The only bring component in the literature; no result published for it \\
\addlinespace
Specialist agent, seen teams \cite{angliss2026} & 0.699--0.740 & Battle policy, not selection; same simulator and protocol \\
\addlinespace
Generalising agent, 72 unseen teams \cite{angliss2026} & 0.564--0.669 & As above; the range in which effects here are found meaningful \\
\addlinespace
Outcome prediction, human data \cite{oblander2025} & 66.43\% acc., 0.730 AUC & Conditions on the six-Pok\'emon team, never on the four brought \\
\addlinespace
Outcome prediction, simulated \cite{reis2024} & 85.4\% acc. & Simplified engine, three-Pok\'emon teams, no human play \\
\addlinespace
Battle-policy ceiling \cite{grigsby2025} & 64--80\% GXE & Gen 1--4 OU singles; no Team Preview exists in the format \\
\addlinespace
Dominance threshold \cite{reis2021} & 55\% & Balance ceiling adopted here as a success floor; see Section~\ref{sec:baseline} \\
\addlinespace
Chance floor (this work) & 1.11\% / 6.67\% & Uniform over the 90 brings, or the 15 lead pairs \cite{angliss2026} \\
\bottomrule
\end{tabular}
\end{table}



\section{Proposed Methodology}
%Problem description, mathematical formulations if needed. Pipeline, experimental protocol, partitions, validation, metrics, and their justification. Baselines and identification and mitigation of risks. Track C specific: formulate MDP, state space, actions space, enviorment transition dynamic, **justified design of the reward function**
\begin{table}[!t]
\centering
\caption{Planned baselines.}
\label{tab:baselines}
\footnotesize
\setlength{\tabcolsep}{4pt}
\renewcommand{\arraystretch}{1.2}
\begin{tabular}{@{}p{2.4cm}p{3.1cm}p{2.0cm}@{}}
\toprule
\textbf{Baseline} & \textbf{Description} & \textbf{Expected level} \\
\midrule
Random & Uniform choice among the 90 legal bring decisions & 1.11\% exact match; 6.67\% on the lead pair alone \\
Usage statistics & Predict the most frequent bring for the observed archetype & Non-trivial; humans weight population usage at 59.3\% \cite{oblander2025} \\
Type heuristic & Greedy type-advantage matching & Reported as sub-optimal \cite{simoes2020} \\
Search builders & Individual Pok\'emon Counter, Minimax Builder (NE via LP), Meta Counter, adapted to the bring decision & Meta Counter strongest in \cite{reis2024} \\
Evolutionary search & GA over the 90 candidates with a learned matchup fitness & Converged in 7--25 generations in \cite{reis2025} \\
Proposed policy & Reinforcement-learning policy over the Team Preview decision & - \\
\bottomrule
\end{tabular}
\end{table}
\subsection{The problem, abstracted}
Running the risk of redundancy but for the sake of clarity: following \cite{angliss2026}, a Pok\'emon VGC battle is a two-player
zero-sum, partially observable, symmetric, stochastic game. Given an open team sheet format, where all the relevant Pokémon specific information (moves, item, stat alignment, ability), is known for both players, take fixed team
configurations $c_1, c_2$ revealed at Team Preview. The game instance is:

\begin{equation}
\mathcal{G}(c_1,c_2) = \langle \mathcal{S}, \{\mathcal{O}_i\}, \{\mathcal{A}_i\}, \mathcal{T}, R, c_1, c_2 \rangle ,
\label{eq:posg}
\end{equation}

\noindent where $\mathcal{S}$ is the set of battle states, $\mathcal{O}_i \subseteq \mathcal{S}$
player $i$'s observation, $\mathcal{A}_i$ the action set for the player (switch out, attack, protect, etc.),
$\mathcal{T}(s' \mid s, a_1, a_2)$ the stochastic transition function (the probability to go into state $s'$ , given the current state $s$ and both players' actions), and
$R(s_T) \in \{-1, 1\}$ is the terminal reward.

We restrict attention to the Team Preview decision that precedes the battle.
Each player selects a bring $b_i$ consisting of an unordered pair of leads and an unordered pair of reserves, drawn from their six.

\begin{equation}
|\mathcal{B}| = \binom{6}{2}\binom{4}{2} = 15 \times 6 = 90 ,
\label{eq:actionspace}
\end{equation}

\noindent giving 90 legal bring decisions per player \cite{angliss2026}.
The learning problem is to find a policy
$\rho(b_1 \mid c_1, c_2)$ maximising the expected terminal reward when the
resulting battle is played out by a shared fixed battle policy, thus removing player skill and
expression from the evaluation.

\subsection{The decision as a Markov Decision Process}
\label{sec:mdp}

We have made clear the learning objective, and now this section makes the
decision process behind it explicit. Fixing the opponent's bring policy and
using a shared battle policy $\pi_b$ for both sides absorbs the opponent into
the environment, reducing the game $\mathcal{G}(c_1,c_2)$ to a single-agent,
finite-horizon Markov Decision Process
$\mathcal{M} = \langle \mathcal{S}_{\mathcal{B}}, \mathcal{A}_{\mathcal{B}}, T_{\mathcal{B}}, R_{\mathcal{B}}, \gamma \rangle$
with a horizon of two steps. The same device is used by \cite{reis2024} and
\cite{reis2025}, which fix a battle policy in order to study the decisions
surrounding it.

\subsubsection{State space}
A state is the information available at Team Preview,
\begin{equation}
s = \langle c_1,\, c_2,\, m,\, \ell \rangle \in \mathcal{S}_{\mathcal{B}},
\end{equation}
\noindent where $c_1$ and $c_2$ are the team configurations of
(\ref{eq:posg}), $m$ is a vector of metagame features derived from usage
statistics over the training window, and $\ell$ is the lead pair already
committed, empty at $t=0$. Under open team sheets the state is fully observed
up to the opponent's concealed stat allocation and their simultaneous bring choice,
the latter of which the fixed opponent policy makes part of the environment.
The MDP state is therefore strictly smaller than the battle state
$\mathcal{S}$: no turn has yet been played.

\subsubsection{Action space}
The 90 alternatives of (\ref{eq:actionspace}) are selected across the two
steps, following the factorisation already present in that expression. At
$t=0$ the agent commits the lead pair,
$a_0 \in \binom{c_1}{2}$, giving 15 alternatives; at $t=1$ it commits the
reserves from the four Pok\'emon that remain,
$a_1 \in \binom{c_1 \setminus a_0}{2}$, giving 6. Their composition is the
bring, $b_1 = (a_0, a_1)$, and
$|\mathcal{A}_{\mathcal{B}}| = 15 \times 6 = |\mathcal{B}| = 90$. Factorising
reduces the per-step branching from 90 to 15 and then 6 without altering the
set of achievable brings.

\subsubsection{Transition dynamics}
The first transition is deterministic and records the commitment,
$T_{\mathcal{B}}(s' \mid s, a_0)$ setting $\ell \leftarrow a_0$ and leaving
$c_1, c_2, m$ unchanged. The second is stochastic and is realised by playing
the battle: both brings are submitted to the simulator and the episode is
executed under $\pi_b$, so that $T_{\mathcal{B}}(\cdot \mid s, a_1)$ is the
composition of the battle transition $\mathcal{T}$ of (\ref{eq:posg}) over the
full episode with the fixed policies of both players. Stochasticity enters
through damage rolls, accuracy checks, secondary effects and speed ties;
\cite{angliss2026} estimates the per-turn branching factor of VGC at up to
$10^{12}$. The agent takes no further action once $a_1$ is committed.

\subsubsection{Reward function}
The reward is sparse and terminal, inheriting $R(s_T)$ from (\ref{eq:posg}):
\begin{equation}
R_{\mathcal{B}}(s, a_0) = 0, \qquad
R_{\mathcal{B}}(s, a_1) = R(s_T) \in \{-1, +1\},
\end{equation}
\noindent with $\gamma = 1$, since the horizon is finite and no intermediate
reward is emitted.

Shaped rewards are common in Pok\'emon because they address credit assignment
over long battles: \cite{simoes2020} decomposes the signal into damage dealt,
faints and damage taken so that the value of a switch is visible before the
battle ends, and \cite{grigsby2025} adds light shaping for damage dealt and
health recovered on top of win and loss. That motivation does not apply here.
The episode is two steps long, so the outcome is never more than two decisions
away and there is no gap for an intermediate signal to bridge.

Shaping would nonetheless introduce a documented risk.  As previously mentioned,
excluding win rate from the fitness function drove a genetic algorithm
toward degenerate parameter values\cite{reis2025}, and \cite{karten2026} document reward
shaping being exploited so that agents bypass intermediate objectives, as well
as negative rewards inducing policy collapse. A sparse outcome reward is not
open to these failures, because it measures exactly the quantity the research
question asks about.

No learned proxy is substituted for $R_{\mathcal{B}}$ in this formulation:
the policy is trained on outcomes recorded in the data. The role of the learned
matchup evaluator introduced in Section~\ref{sec:evaluator} is confined to the
evolutionary baseline, for reasons given there.

\subsection{Observability of the bring in recorded play}
\label{sec:observability}

The formulation above assumes that a logged battle supplies the bring each
player selected. The issue raised in Section~\ref{sec:rltechniques} is that it does
not always do so, and the design that follows depends on knowing exactly how
often and in what way. The Regulation~M-A log set described in
Section~\ref{sec:axisc} was therefore parsed directly and measured before any
model was specified, over the whole of its 6\,905 battles and therefore all
13\,810 of its player-brings.

What a log records is not a bring but a sequence of appearances. The two
Pok\'emon a player sends out immediately after Team Preview are written to the
log as switch-ins at the start of the battle, so \textbf{the lead pair is
recovered exactly in every battle}. A reserve, by contrast, appears in the log
only if it is brought onto the field at some point, and a battle can end before
that happens: the winning player has no reason to switch once the position is
decided, and the losing player may forfeit while a member is still on the
bench. Formally, for a logged bring the observation is

\begin{equation}
o = \langle L,\, V \rangle, \qquad L \subset c,\ |L| = 2, \quad V \subseteq c \setminus L ,
\label{eq:obs}
\end{equation}

\noindent where $c$ is that player's six-Pok\'emon team configuration, $L$ the
observed lead pair and $V$ the set of reserves that were revealed,
$v = |V| \in \{0,1,2\}$. The distribution of $v$ over the
13\,810 player-brings measured is 84.1\% with $v=2$ (the complete bring is
recorded), 12.0\% with $v=1$, and 3.9\% with $v=0$. Both brings of a battle are
complete in 72.5\% of cases.

Two properties of this missingness matter more than its rate. The first is that
it is confined to the second step of the MDP: $a_0$ is supervised in 100\% of
recorded player-brings and only $a_1$ is ever ambiguous. The second is that it
is not random. Complete reveals are far more frequent among losing players
(92.8\%) than among winning ones (73.2\%), for the obvious reason that a player
who is winning is not forced to commit their remaining Pok\'emon. The rate of
disclosure is therefore correlated with the outcome, which has a direct
consequence for the design: discarding incomplete records would not
produce a smaller unbiased sample but a sample enriched in losing brings, and
the number of revealed Pok\'emon must never be supplied to any model as a
feature, since it is a proxy for the label the model is meant to predict.

\subsection{Predicting the complete bring}
\label{sec:whyfull}

An obvious response to Section~\ref{sec:observability} would be to reduce the
problem to the part of it that is always observed, predicting only the lead
pair and its 15 alternatives. That reduction is rejected here, for five
reasons.

\paragraph{The decision the game poses is the bring} A player at Team Preview
submits four Pok\'emon at once. The reserves are committed in the same instant
as the leads and cannot be revised afterwards. A model that predicts leads
answers a question no player is ever asked, and the research question of this
work is stated over the decision as the format defines it, not over the portion
of it that happens to be easiest to observe.

\paragraph{The reserves determine what the leads are worth} Whether a lead pair
is a good choice depends on what is available to come in behind it. Switching
is a central mechanic of the format, and \cite{simoes2020} report that the
value of switching is one of the behaviours an agent acquires without any prior
domain knowledge, which is to say it is a real and learnable component of
competitive value rather than an incidental one. A lead pair scored without its
reserves is scored against an unspecified continuation.

\paragraph{Comparability with the existing work} Ninety is the size
this decision has been given in the only published treatment of it
\cite{angliss2026}, and the Selection Phase of \cite{reis2021} also selects
a sub-team rather than a front line. Reducing the space to 15 would place the
results of this work outside comparison with either, and would understate the
difficulty of the task by a factor of six.

\paragraph{The labels are partial, not absent} The distinction is what allows us to still consider the whole bring, even if some of it is unseen. In each incomplete record it is known exactly how many reserves
are missing and exactly which Pokémon they could be, so what the log supplies
is a set of brings consistent with what was seen,

\begin{equation}
\mathcal{C}(o) = \{\, b = (L, R) \;:\; V \subseteq R \subseteq c \setminus L \,\},\\
\qquad |\mathcal{C}(o)| = \binom{4-v}{2-v} ,
\label{eq:consistent}
\end{equation}

\noindent where $R$ is the reserved pair, so that $|\mathcal{C}(o)| = 1$ when $v=2$, $3$ when $v=1$ and $6$
when $v=0$. Weighted by the frequencies of Section~\ref{sec:observability}, the
average record narrows the 90 alternatives to 1.44 of them. This is a
supervision problem and not a reason to change the question. The corresponding
objective replaces the log-likelihood of a single label with the log-likelihood
of the consistent set,

\begin{equation}
\mathcal{L} \;=\; -\sum_{o \,\in\, \mathcal{D}} \log \!\!\sum_{b \,\in\, \mathcal{C}(o)} \rho(b \mid s) ,
\label{eq:partial}
\end{equation}

\noindent where $\mathcal{D}$ is the set of logged brings in the training
window, and which reduces exactly to ordinary cross-entropy on the 84.1\% of
records where $\mathcal{C}(o)$ is a singleton, and which spreads likelihood over
the consistent completions. Since $\mathcal{C}(o)$ always contains the true
bring, the objective is never trained against a target known to be wrong.

\paragraph{Evaluation is unaffected} The limitation lives in the supervision
and nowhere else. The agent emits a complete bring at every decision, that
bring is submitted in full to the simulator, and the primary measure of this
work is the win rate of the battles that result. No quantity in the answer to
the research question is read off a log at all, so no part of it inherits the
ambiguity.

To this can be added that reducing the problem to the lead pair would collapse
the two-step Markov Decision Process of Section~\ref{sec:mdp} to a single step,
discarding the factorisation of (\ref{eq:actionspace}) and, with it, the only
stage of the decision at which the counter-selection reasoning described by
\cite{crane2021} operates: in VGC, the reserve pair is precisely where a player holds
their response to what the opponent is expected to lead with. Simply put, this makes the problem more interesting and the solution an actual novelty.

%  optional, if you keep the learned matchup evaluator 
\subsection{Learned Matchup Evaluator}
\label{sec:evaluator}
Searching the space of bring decisions requires a way to score a candidate
without playing it. The direct alternative, simulating battles for each of
the 90 candidates, is impractical. Pok\'emon is highly stochastic, and
\cite{galway2008} notes that stochastic game environments produce noisy fitness
evaluations requiring a large number of games per candidate before the estimate
is reliable. Multiplying 90 candidates by the number of battles needed to
distinguish them, and again by the number of search iterations, places
simulation-based scoring beyond the computational capabilities of this work. We
therefore adopt a learned matchup evaluator as the objective that search
maximises, similar to the approach of \cite{reis2024}, in which a neural network
trained to predict match outcomes replaces battle simulation as the criterion
for team search.

Since VGC is symmetric and zero-sum, it can then be collapsed into the following: Let $f$ map a bring to a $K$-dimensional embedding and let $M$ be an
antisymmetric matrix. We adopt the bilinear matchup formulation of
\cite{oblander2025}, which models the win probability between two
six-Pok\'emon teams, and apply it at the level of the bring. The estimated
advantage of $b_1$ against $b_2$ is
\begin{equation}
\hat{v}(b_1, b_2) = \sigma\!\left( f(b_1)^{\top} M f(b_2) \right), \qquad M = -M^{\top},
\label{eq:bilinear}
\end{equation}
\noindent where $\sigma$ denotes the logistic function. Since $M$ is
antisymmetric, exchanging the arguments negates the bilinear score, and
$\sigma(-x) = 1 - \sigma(x)$ therefore gives
$\hat{v}(b_1,b_2) = 1 - \hat{v}(b_2,b_1)$ by construction.

The embedding function $f$ is instantiated as follows. Each Pok\'emon $x$ is
mapped to a feature vector $\phi$ by an embedding layer over its species,
moves, held item, and ability. Given a bring $b$ with lead pair $L$
and reserve pair $R$,
\begin{align}
h(b) &= \mathrm{ReLU}\Big(\beta_h + W_L \sum_{x \in L} \phi_x + W_R \sum_{x \in R} \phi_x\Big), \\
f(b) &= \sigma\big(\beta_z + W_z\, h(b)\big),
\end{align}
\noindent where $\beta_h, \beta_z$ are biases and $W_L, W_R, W_z$ learned
projections. Summing within each subset makes $h$ invariant to the order in
which the two leads or the two reserves are listed, since using distinct
projections $W_L \neq W_R$ prevents the two roles from being interchangeable.
This is the structural counterpart of the action space defined in
(\ref{eq:actionspace}), in which the 90 alternatives are counted without regard
to order within either pair.

This functional form is chosen for three reasons. First, the antisymmetry
constraint enforces the zero-sum property of the game structurally rather than
requiring the model to approximate it from data. Second, the constraint
is obtained at no cost in predictive accuracy: \cite{oblander2025} reports
66.43\% test accuracy for the bilinear specification against 66.40\% for an
unconstrained matchup function on the same data, while
outperforming pre-trained embeddings. Third, the bilinear form
places the matchup relation in a linear structure over the embedding space,
so that the geometry of $f$ can be inspected to characterise which
configurations counter which, an added layer of interpretability that a general
non-linear predictor does not provide.

To arrive at our \textbf{fitness function} $\hat{v}$ starting from \cite{oblander2025}'s, three adaptations are required. The original formulation embeds an unordered set of six Pok\'emon and
accordingly constrains $f$ to be permutation-symmetric across all slots. A
bring is not fully unordered: the two leads occupy the field at the start of
the battle while the two reserves do not. We therefore retain permutation
symmetry within each pair while distinguishing the pairs themselves, computing
separate order-invariant aggregations over the lead and reserve subsets and
combining them through distinct projections. Once again, this is how it preserves the invariance
implied by the action space, in which the $\binom{6}{2}\binom{4}{2} = 90$
alternatives are counted without regard to order within either pair, while
allowing the model to represent the distinction between leading with a
Pok\'emon and holding it in reserve. 

The original
also includes an additive player-skill term so that $M$ captures matchup
structure residual of player ability; in our evaluation both sides are played
by an identical fixed battle policy, so no such correction is required.

The third adaptation follows from Section~\ref{sec:observability}. $\hat{v}$ is
trained on human play rather than on simulated matches: each Open Team Sheet
battle log in \cite{angliss2026} yields both players' six-Pok\'emon teams, the
observations $o_1, o_2$ of their two brings, and the outcome. In the 72.5\% of
battles where both brings are recorded in full this is a supervised target of
the form $(b_1, b_2, \text{winner})$ exactly as intended. Where it is not, the
training example is a pair of consistent sets rather than a pair of brings, and
the likelihood of the observed outcome is marginalised over their product in
the manner of (\ref{eq:partial}),
\begin{equation}
\Pr(w{=}1 \mid o_1, o_2) = \!\!\!\sum_{b_1 \in \mathcal{C}(o_1)} \sum_{b_2 \in \mathcal{C}(o_2)} \!\!\! q(b_1)\, q(b_2)\, \hat{v}(b_1,b_2) ,
\label{eq:marginal}
\end{equation}
\noindent with $q$ uniform over each consistent set, the complementary
expression for $w{=}2$, and at most $6 \times 6 = 36$ terms in the double sum.
No battle is discarded, which for the reasons given in
Section~\ref{sec:observability} matters: the discarded 27.5\% would not be a
random subset.

\subsection{Why both $\hat{v}$ and $R_{\mathcal{B}}$?}

The reward $R_{\mathcal{B}}$ of Section~\ref{sec:mdp} is available only for
battles that were actually played. For training the selection policy this poses
no difficulty: each logged battle is a complete episode, supplying the state,
the bring and the outcome directly. For searching it gets a little trickier. A procedure that
explores the 90 alternatives generates candidates no player selected and for
which no outcome was ever recorded, and ranking them requires estimating
$R_{\mathcal{B}}$ for pairs absent from the data.

The learned evaluator provides that estimate. Where $R_{\mathcal{B}} \in
\{-1,+1\}$ is a realised outcome, $\hat{v}(b_1,b_2) \in (0,1)$ approximates the
probability that $b_1$ defeats $b_2$, so that
\begin{equation}
    \hat{v} \approx \frac{\mathbb{E}[R_{\mathcal{B}}] + 1}{2}
\end{equation}
where $\mathbb{E}$ denotes expectation over the stochasticity under the fixed battle
policy $\pi_b$. The two are the same quantity observed at different levels: one
discrete and sampled, the other continuous and predicted.
The role of $\hat{v}$ is confined to search. It is an approximation, learned
from data, used to rank candidate brings during training at a cost that
simulation cannot match. 

To be clear, $\hat{v}$ is not used to answer the research question. That
answer comes from played battles: the selected bring is executed under a fixed
battle policy shared by both sides, over 1000 battles per team pair with the
sides exchanged, and measured against a population-usage baseline. Evaluating
the policy with the same function it was trained to maximise would be
circular, and the separation is maintained deliberately: $\hat{v}$ measures
whether a bring \emph{should} win, and the evaluation measures whether it
\emph{does}. It is worth noting that this separation is what keeps the partial
supervision of Section~\ref{sec:observability} out of the result: the
ambiguity enters $\hat{v}$ and the selection policy, both of which are trained
from logs, and never enters the win rate, which is produced by simulation.

One consequence of training on human logs should be noted. $\hat{v}$ learns
the regularities of play between human opponents, whereas evaluation is
conducted under a fixed battle policy. If that policy is a reasonable proxy for
competent human play, the learned advantages should transfer; if it is not, a
divergence between $\hat{v}$ and the observed win rate is itself an
informative result, and is reported as such.

\subsection{Experimental Pipeline}

It comprises four stages.

\subsubsection{Data extraction} Battle logs are obtained from the Open Team Sheet
VGC dataset released with \cite{angliss2026} and parsed directly from the raw
Showdown protocol rather than through the accompanying log-to-trajectory
tooling, which presupposes that both brings are known. From each log the six
Pok\'emon of each side are read from the team-sheet lines, the lead pair of
each side from the two switch-ins recorded immediately after the start of the
battle, the revealed reserves from every subsequent entry onto the field, and
the winner from the terminal line, matched to a side by the player
declarations. Nicknames, form changes and mid-battle transformations are
normalised to the base species before matching, since a Pok\'emon may appear
under a different identifier once it changes form. Each battle therefore yields
$(c_1, c_2, o_1, o_2, w)$, with $o_i$ as defined in (\ref{eq:obs}), which
collapses to the intended $(c_1, c_2, b_1, b_2, w)$ whenever both observations
are complete.

\subsubsection{Baseline construction} Bring frequencies conditioned on team archetype
are computed over the training window to instantiate the usage-statistics
baseline. Since a frequency table is itself a label-consuming estimate, incomplete
records contribute fractional counts spread uniformly over $\mathcal{C}(o)$
rather than being dropped, for the reason given in Section~\ref{sec:observability}.

\subsubsection{Model training} The matchup evaluator $\hat{v}$ is fitted to the
extracted records under (\ref{eq:marginal}), and the selection policy is
trained over the MDP of Section~\ref{sec:mdp} under (\ref{eq:partial}). No
feature derived from the length of the battle or from the number of Pok\'emon
revealed is made available to either model, i.e. winning while revealing less Pokémon is not rewarded more.

\subsubsection{Evaluation} Selected brings are played out on a locally hosted
Showdown, a battle simulator, server under a fixed battle policy, and win rates are compared against
the baselines.

\subsection{Experimental Protocol}
\label{sec:protocol}

\subsubsection{Partitions} Battles are split \textbf{temporally} rather than at random:
the training set precedes the validation and test sets in time. The Pok\'emon
metagame evolves continuously \cite{karten2026}, and a random split would allow
a model to exploit knowledge of later trends when predicting earlier battles.
The usage statistics defining the baseline are computed on the training window
only, for the same reason.

\subsubsection{Validation} Model selection uses a held-out set of team pairs unseen
during training, following the generalization protocol of \cite{angliss2026},
which reserves unseen teams to separate performance from memorisation.

\subsubsection{Metrics} Four are reported, of which one is primary and three are
diagnostic. \emph{Downstream win rate} is the primary measure, as it is the
quantity named in the research question, and it is obtained entirely from
simulated battles and therefore carries none of the label ambiguity of
Section~\ref{sec:observability}.

The three diagnostics answer \emph{why} a win rate is or is not attained, and
each is paired with the portion of the data on which it is well defined.
\emph{Lead-pair exact match} measures whether the opponent's lead pair is
recovered, over 15 alternatives and against a chance floor of 6.67\%; it is
defined on 100\% of the test battles. \emph{Per-slot recall} measures how many
of the revealed members of the opponent's bring are identified, degrades
gracefully where exact match does not, and is likewise defined on every battle,
since it asks only about Pok\'emon that were seen. \emph{Full-bring exact
match}, over 90 alternatives and against a chance floor of 1.11\%, is defined
only where the opponent's bring was recorded in full. It is reported in two
forms: a strict figure computed on the complete-case subset, stated together
with the caveat that this subset over-represents losing players
(Section~\ref{sec:observability}) and is therefore not a random sample of test
battles; and a permissive figure counting a prediction as correct when it lies
in $\mathcal{C}(o)$, which is unbiased in coverage but optimistic in value. The
two bracket the quantity of interest, and both are given rather than either
alone.

\subsubsection{Sample sizes and reporting} Each matchup is evaluated over 1000 battles
with the sides exchanged, matching the protocol of \cite{angliss2026}, which
reports an error of approximately $\pm 0.03$ at that sample size; the swap
follows \cite{reis2021}, where teams are exchanged between players so that
results are not confounded by team strength. Results are reported as
distributions over at least ten random seeds and against a population of
opponents rather than a single fixed one, as argued by \cite{liwellman2024}.

\subsection{Baselines}

Table~\ref{tab:baselines} lists the planned baselines. Of them, the
usage-statistics predictor is the one the proposed policy must beat:
\cite{oblander2025} measures human players placing 59.3\% of their decision
weight on the population action distribution, so a model reproducing that
distribution reproduces most of human behaviour and is a demanding rather than
a trivial comparison. Every baseline is fitted on the same records and under
the same treatment of incomplete ones as the proposed policy, so that no part
of a reported difference is attributable to one method having been given
cleaner data than another.

\subsection{Risks and Mitigation}
A localized summary of the decisions previously justified that are looking to minimize misrepresentation of the solution.

\subsubsection{Data extraction} This risk is now discharged rather than
projected. The bring was confirmed recoverable from the released logs by
parsing 6\,905 of them directly, and the recovery is exact for the lead pair in
every battle and for the full bring in 84.1\% of player-brings
(Section~\ref{sec:observability}); the reconstruction procedure of
\cite{grigsby2025} is therefore not required. Two of the 13\,810 player-brings
parsed produced an internally inconsistent record, attributing more than four
Pok\'emon to a four-slot bring, which indicates a residual edge case in the
handling of form changes or of the illusion mechanic. These represent 0.014\%
of records and are excluded; the parser is re-verified against a manually
annotated sample before the full extraction is used.

\subsubsection{Partial supervision} The marginalised objectives
(\ref{eq:partial}) and (\ref{eq:marginal}) assume that, conditional on the
state, which reserves were revealed carries no information about which bring
was chosen. That assumption is imperfect, since disclosure depends on how the
battle developed. It is mitigated in three ways: no feature derived from the
number of revealed Pok\'emon is given to any model, so the correlation cannot
be exploited; the metrics affected are reported on both the complete-case and
the consistent-set convention, which bound the effect from either side; and the
primary result is a simulated win rate, which does not depend on the labels at
all.

\subsubsection{Complete-case bias} Filtering to fully revealed brings would be
the simplest response to the limitation and is explicitly not taken. Complete
reveals occur in 92.8\% of losing brings against 73.2\% of winning ones, so the
filtered sample is enriched in brings that lost. Retaining every record under
(\ref{eq:partial}) avoids importing that bias into the training distribution.

\subsubsection{Reward specification} Addressed by the sparse outcome reward argued for
above, with any shaping term ablated before it is retained.

\subsubsection{Noisy fitness} Scoring candidates by simulation is unreliable in a
stochastic game \cite{galway2008}; mitigated by the learned evaluator, with a
fixed battle budget per candidate and reported variance wherever simulation is
used.

\subsubsection{Distribution shift} $\hat{v}$ is trained on human play but evaluated
under a fixed battle policy. Divergence between predicted and observed
advantage is measured and reported rather than assumed away.

\subsubsection{Overfitting to fixed teams} \cite{huang2019} report a specialised agent
falling to 77 wins in 500 games outside its training teams, and
\cite{angliss2026} find single-team agents generalising worst. Mitigated by
reporting results on held-out team pairs alongside seen ones.

\subsubsection{Exploitability} Agents in this domain are close to fully exploitable
\cite{angliss2026}. A best-response win rate is reported so that robustness is
stated rather than implied.

\subsubsection{Scope} As it's been made explicit, no battle policy is trained; existing agents are reused, as in
\cite{reis2024} and \cite{reis2025}.


\begin{thebibliography}{16}

\bibitem{hu2024}
C.~Hu, Y.~Zhao, Z.~Wang, H.~Du, and J.~Liu, ``Games for artificial intelligence research: a review and perspectives,''
\emph{IEEE Transactions on Artificial Intelligence}, vol.~5, no.~12, pp.~5949--5968, Dec.~2024.

\bibitem{vgchandbook}
The Pok\'emon Company International, ``Play! Pok\'emon Video Game Championships
tournament handbook,'' 2026. [Online].
Available: \url{https://www.pokemon.com/static-assets/content-assets/cms2/pdf/play-pokemon/rules/play-pokemon-vgc-tournament-handbook-en.pdf}
[Accessed: Sep. 3, 2026].

\bibitem{karten2026}
S.~Karten \emph{et al.}, ``The Pok\'eAgent challenge: competitive and long-context learning at scale,''
in \emph{Advances in Neural Information Processing Systems 38 (NeurIPS 2025) Competition Track}, 2025.

\bibitem{angliss2026}
C.~L. Angliss, J.~Cui, J.~Hu, A.~Rahman, and P.~Stone, ``VGC-Bench: towards mastering diverse team strategies in competitive Pok\'emon,''
in \emph{Proc. 25th Int. Conf. Autonomous Agents and Multiagent Systems (AAMAS)}, Paphos, Cyprus, 2026.



\bibitem{reis2021}
S.~Reis, L.~P. Reis, and N.~Lau, ``VGC AI competition - a new model of meta-game balance AI competition,''
in \emph{2021 IEEE Conference on Games (CoG)}, 2021, pp.~1--8.

\bibitem{simoes2020}
D.~Sim\~oes, S.~Reis, N.~Lau, and L.~P. Reis, ``Competitive deep reinforcement learning over a Pok\'emon battling simulator,''
in \emph{2020 IEEE International Conference on Autonomous Robot Systems and Competitions (ICARSC)}, Ponta Delgada, Portugal, 2020, pp.~40--45.

\bibitem{grigsby2025}
J.~Grigsby, Y.~Xie, J.~Sasek, S.~Zheng, and Y.~Zhu, ``Human-level competitive Pok\'emon via scalable offline reinforcement learning with transformers,''
\emph{Reinforcement Learning Journal}, 2025.

\bibitem{huang2019}
D.~Huang and S.~Lee, ``A self-play policy optimization approach to battling Pok\'emon,''
in \emph{2019 IEEE Conference on Games (CoG)}, London, U.K., 2019, pp.~1--4.

\bibitem{reis2024}
S.~Reis, R.~Novais, L.~P. Reis, and N.~Lau, ``An adversarial approach for automated Pok\'emon team building and metagame balance,''
\emph{IEEE Transactions on Games}, vol.~16, no.~2, pp.~365--375, Jun.~2024.

\bibitem{reis2025}
S.~Reis, A.~L. Martins, R.~Novais, and F.~Alves, ``A new rules balance track for the Pok\'emon VGC AI competition 2.0,''
in \emph{2025 IEEE Conference on Games (CoG)}, Lisbon, Portugal, 2025, pp.~1--8.

\bibitem{rodriguez2024}
G.~Rodriguez, E.~Villanueva, and J.~Bald\'eon, ``Enhancing Pok\'emon VGC player performance: intelligent agents through deep reinforcement learning and neuroevolution,''
in \emph{HCI in Games}, ser. Lecture Notes in Computer Science, vol.~14730. Cham, Switzerland: Springer, 2024, pp.~275--294.

\bibitem{li2025}
Z.~Li, Q.~Ji, X.~Ling, and Q.~Liu, ``A comprehensive review of multiagent reinforcement learning in video games,''
\emph{IEEE Transactions on Games}, vol.~17, no.~4, pp.~873--892, 2025.

\bibitem{galway2008}
L.~Galway, D.~Charles, and M.~Black, ``Machine learning in digital games: a survey,''
\emph{Artificial Intelligence Review}, vol.~29, no.~2--4, pp.~123--161, 2008.

\bibitem{liwellman2024}
Z.~Li and M.~P. Wellman, ``A meta-game evaluation framework for deep multiagent reinforcement learning,''
in \emph{Proc. 33rd Int. Joint Conf. Artificial Intelligence (IJCAI)}, 2024, pp.~148--156.

\bibitem{oblander2025}
S.~Oblander, ``Representation learning for behavioral analysis of complex competitive decisions,''
Social Science Research Network, Rochester, NY, Working Paper, 2025. [Online]. Available: \url{https://papers.ssrn.com/abstract=4969862}

\bibitem{crane2021}
D.~Crane, Z.~Holmes, T.~T. Kosiara, M.~Nickels, and M.~Spradling,
``Team counter-selection games,'' in \emph{Proc. IEEE Conf. Games (CoG)}, 2021, pp.~1--8.

\end{thebibliography}
\vspace{12pt}

\end{document}
